\section{What remains when a boundary encloses an event}
\label{sec:eventclosure}

An estimate on a late regular tube does not determine the flux on a
complete boundary starting at time zero. We now give a quantitative
description of this difference. It also shows why the extra terms in
\eqref{eq:organizing} cannot be suppressed by taking a thinner strip.

Let a fixed core be
\[
 W_c(Y)=[0,a]+i[-Y,Y]\;\times\;\{|z-x_*|\le R\},
 \qquad x_*\in\R,\quad 0<Y\le Y_0.
\]
Assume real type, a nonzero jet on its boundary for these $Y$, and
that the only joint zeros in the core are a fixed finite set of real
events in its interior. Denote their total complex charge by $Q$.
The initial face, the whole face $t=a$, and the spatial wall are
uniformly separated from the joint zeros. Suppose on their union
for $|s|\le Y_0$ that
\begin{equation}
 |G|\ge b>0,\qquad \|DG\|_{\rm op}\le D.
 \label{eq:corefloor}
\end{equation}
These are compact boundary quantities, not assumptions about a floor
at the events. Assume $m$ disjoint real root tubes attach along disks
of the $t=a$ face and continue to the terminal cap. The complement
of those disks remains part of the boundary.

For the canonical flow, a fixed finite set of real simple branches
that remains simple for $t\ge a$ admits such tubes for every finite
$T$. A compact initial segment has a common positive complex-time
width and inverse radius by the implicit function theorem. The
late part is covered by Section~\ref{sec:source}. Combining them,
and decreasing the common width if necessary, permits a fixed small
$Y>0$ and radius $r/L_T$ for all sufficiently large $T$. For
example, branches starting at $a>1/2$ stay simple: de Bruijn's
reality theorem and the local splitting lemma rule out a later
collision. This continuation statement neither asserts an actual
earlier theta event nor supplies a height-uniform core floor.

\begin{theorem}[Quantitative closure allocation]
\label{thm:core}
In this fixed-core geometry let $\mathcal R_T^{\rm all}$ be the
signed sum of the initial, joining and local spatial faces, including
the initial sign $-\II_0$. Set
\begin{equation}
 B=\frac{D^3}{2\pi^2b^3}(4\pi R^2+4\pi Ra).
 \label{eq:Bcore}
\end{equation}
Every projective tube wall has zero flux, and
\begin{equation}
 |\mathcal R_T^{\rm all}|\le BY,\qquad
 \boxed{\displaystyle
 \JJ_+=\JJ_-=
 \frac{Q-\II_{T,m}-\mathcal R_T^{\rm all}}2.}
 \label{eq:coreidentity}
\end{equation}
For the following asymptotic assertions, assume that $H$ is the
canonical flow \eqref{eq:theta} and that the attached roots have
fixed positive ranks. If $Y_T=S/L_T^2$ and $\epsilon_T=r/L_T$, with
fixed $S,r$ and fixed attached population, then
\begin{equation}
 \II_{T,m}=O(mL_T^{-2}),\qquad
 \left|\JJ_\pm-\frac Q2\right|
 \le\frac{BY_T+|\II_{T,m}|}{2}
 =O((B+m)L_T^{-2}).
 \label{eq:halfcharge}
\end{equation}
For fixed admissible $Y>0$, let $R_c(Y)$ be the signed sum of the
three non-imaginary faces of the unmodified core. Then
\begin{align}
 \II_{T,m}&\longrightarrow \frac{mr^2Y}{\pi},\nonumber\\
 \mathcal R_T^{\rm all}&\longrightarrow R_c(Y),
       &|R_c(Y)|&\le BY,\nonumber\\
 \JJ_\pm&\longrightarrow
      \frac{Q-R_c(Y)}2-\frac{mr^2Y}{2\pi}.
 \label{eq:fixedwidthcore}
\end{align}
The errors in these fixed-width limits are
$O(L_T^{-2})+O(\varepsilon_T/L_T^2)$ with fixed core and
population constants.
\end{theorem}
\begin{proof}
The normalized sphere-volume expression for \eqref{eq:forms}
gives a comass bound $D^3/(2\pi^2b^3)$ on the faces in
\eqref{eq:corefloor}. Their total three-volume is at most
$Y(4\pi R^2+4\pi Ra)$. Removing the attachment disks only
reduces this upper bound. This proves the first inequality.

Simultaneous conjugation $(t,s,x,y)\mapsto(t,-s,x,-y)$ preserves
the ambient orientation, exchanges the outward imaginary faces,
and sends $\Aconn$ and $d\Aconn$ to their negatives. It therefore
preserves $\CC$ and proves $\JJ_+=\JJ_-$. The complete boundary
has charge $Q$ by Lemma~\ref{lem:degree}, yielding
\eqref{eq:coreidentity}. Equation~\eqref{eq:shrinkingcap} proves
\eqref{eq:halfcharge}.

For fixed width, the removed attachment caps at $t=a$ have size
$O(k_{\epsilon_T})=O(L_T^{-2})$, because their center derivatives
are nonvanishing on a fixed compact segment and their lifted phase
increments are bounded. Thus
$\mathcal R_T^{\rm all}=R_c(Y)+O(L_T^{-2})$.
The terminal estimate \eqref{eq:capheadline} and the exact balance
give \eqref{eq:fixedwidthcore} with the stated errors.
\end{proof}

The half-charge limit is a characterization of the remaining flux,
not an independent upper bound on the event count. Its constants
belong to a fixed core. Allowing height, multiplicity or the number
of cores to grow requires new estimates for $b,D,R,a$; none is
implicit here. Shrinking $Y$ also changes the initial cap being
integrated. At fixed $Y$, the initial core and its initial integral
remain fixed, and the nonzero limits in \eqref{eq:fixedwidthcore}
give the appropriate comparison.

\subsection{Two exact controls}

The quadratic heat solution
$H(\tau,z)=z^2-2\tau$ has charge one at the origin. Its inverse
root disks are
\[
 Z_\pm(\tau,q)=q\pm\sqrt{q^2+2\tau}.
\]
For $t\ge a>0$ and $\epsilon^2<2a$, these are disjoint regular
tubes on a sufficiently small imaginary-time strip. Their two-root
terminal cap is exactly
\begin{equation}
 \II_T=\frac{k_\epsilon}{\pi}\arctan(Y/T).
 \label{eq:quadraticcap}
\end{equation}
Indeed $\DD=-8\tau$, so the vertical phase increment is
$2\arctan(Y/T)$. The cap tends to zero and the sides have zero
flux, although an attached core still has charge one. The two
inverse branches pinch at $q^2=-2\tau$ when continued backward;
one cannot extend a fixed-radius regular disk bundle through the
event. This example checks the need for a joining region.

For $\lambda>0$, the exact heat solution
$H(\tau,z)=e^{\lambda^2\tau}\sin(\lambda z)/\lambda$
has stationary simple roots. On any fixed finite collection of
their projective disks, $\Delta$ is constant but
$\DD$ has phase rate $m\lambda^2$. Hence
\begin{equation}
 \II_T=\frac{k_\epsilon m\lambda^2Y}{\pi}.
 \label{eq:sinecap}
\end{equation}
This cap is nonzero; the full regular tube still has zero charge.
It checks that root positions alone do not determine the physical
cap. The leading phase in Theorem~\ref{thm:headline} is the
canonical, increasing-frequency counterpart of this calculation.

Both controls satisfy the heat equation but lack the unchanged
theta source. Conversely, the source proof uses only the displayed
positive density bounds and concentration: it is not claimed to
distinguish theta from every other source with the same asymptotic
properties. No collision exclusion can be inferred merely from
the late clock pattern.

\section{What the formula now makes accessible}
\label{sec:outlook}

The initial-to-future identity \eqref{eq:organizing} and its
independent-bound consumer \eqref{eq:eventbound} are the organizing
framework. The estimates proved here supply a terminal cap, its
Vandermonde and analytic-factor decomposition, exactly vanishing
projective sides, and relative late-edge asymptotics for a growing
population. They also identify a substantial surviving term at a
fixed event core. These are concrete bounds with different domains
of validity; their domains must match before they can be combined
to bound events.

The next estimate needed for a full time-zero application is on the
complete earlier attachment and imaginary-time boundary, with
constants controlled as the spatial height grows. A meaningful
first target is a bound for that contribution divided by the
enclosed zero population. It need not be small enough to exclude
every event. One could then use an independently obtained lower
bound on the fixed initial cap in \eqref{eq:eventbound}. Neither
that lower bound nor a uniform early-boundary estimate is proved
here. The fixed-core half-charge formula cannot replace them,
because it already contains the unknown charge.

There is also a separate real-variable route. At fixed positive
heat time, the high-height estimates of Polymath yield zero-free
moving barriers and finite budgets for later real births. Those
budgets do not estimate the physical imaginary-time integrals in
this paper. Likewise a pair-correlation law for real zeros does
not supply the analytic factor in \eqref{eq:organizing}. Keeping
these distinctions explicit leaves a modest but useful conclusion:
the topological formula has quantitatively tractable pieces, and
the remaining task is located on a specified boundary rather than
hidden in a presumed vanishing cap.

\section*{Acknowledgements and research status}
This manuscript was drafted for Edward Baker with substantial
assistance from OpenAI's Codex, recorded model \texttt{gpt-6-astra},
reasoning effort \emph{max}. LLM assistance was used for derivation,
literature comparison, checking and exposition. The checks in this
round are by the same assistant; they are not independent specialist
or human review. The general transport and fixed-population cap
arguments consolidate earlier LLM-assisted research notes in the
author's archive~\cite{Archive}. The present continuation adds the
narrow-strip growing-population estimates, the linear-population
Vandermonde logarithm bound and the proportional-interval edge
asymptotic. The companion working manuscript on zero counts and
localized Hodge energy remains a separate catalogue.

All asymptotic statements have analytic proofs. No numerical onset,
theta-zero certificate, priority certification, proof of RH, or
proof of initial simplicity is asserted. The accompanying literature
and proof review records source access, overlap, control examples
and the remaining novelty question.
